\section{One efficient operating frontier, two questions}
\label{sec:frontier}

Choose any nonzero direction \(\vect d\) and write
\(\gls{sym:i}=\sqrt{\gls{sym:rho}}\,\vect d\), where
\(\gls{sym:rho}\geq0\) is a squared radial scale. Define
\begin{equation}
 S(\vect d)=\vect d^{\mathsf T}\gls{sym:rmat}\vect d,\quad
 T(\vect d)=\vect d^{\mathsf T}\gls{sym:hmat}\vect d,\quad
 V(\vect d)=\vect d^{\mathsf T}\gls{sym:qmat}\vect d.
\end{equation}
Homogeneity gives
\begin{equation}
 \gls{sym:pcu}=\gls{sym:rho}S,\qquad
 \gls{sym:torque}=\gls{sym:rho}T,\qquad
 U^2=\gls{sym:rho}V.
 \label{eq:radial-separation}
\end{equation}
Along that ray, the two physical resources permit
\begin{equation}
 \rho_P=\frac{\gls{sym:pcumax}}{S},\qquad
 \rho_U=\begin{cases}\gls{sym:umax}^2/V,&V>0,\\+\infty,&V=0,\end{cases}
 \qquad \rho_{\max}=\min(\rho_P,\rho_U).
 \label{eq:radial-limits}
\end{equation}
Imagine walking outward from the origin along one current direction. The first
surface reached---loss ellipsoid or voltage cylinder---sets the admissible
radius. Maximum positive torque is therefore
\begin{equation}
 \gls{sym:mmax}=\max_{\vect d\neq0}T(\vect d)
 \min\!\left(\frac{\gls{sym:pcumax}}{S(\vect d)},
 \frac{\gls{sym:umax}^2}{V(\vect d)}\right),
 \label{eq:mmax-direction}
\end{equation}
with the voltage fraction interpreted as infinite for a voltage-null
direction.

For a requested positive torque, a direction with \(T>0\) requires
\begin{equation}
 \rho_M=\frac{\gls{sym:mref}}{T},\qquad
 P_{\mathrm{dir}}=\gls{sym:mref}\frac{S}{T},\qquad
 U_{\mathrm{dir}}^2=\gls{sym:mref}\frac{V}{T}.
 \label{eq:minloss-direction}
\end{equation}
Minimum loss chooses the direction with smallest \(S/T\) among those satisfying
\(M_{\mathrm{ref}}V/T\leq U_{\max}^2\). Equations
\cref{eq:mmax-direction,eq:minloss-direction} are the same direction search;
only the radial question differs.

Define the efficient frontier
\begin{equation}
 \gls{sym:pmin}=\min\{P_{\mathrm{Cu}}(\vect i):M(\vect i)=M,
 U^2(\vect i)\leq U_{\max}^2\}.
 \label{eq:pmin-frontier}
\end{equation}
For \(0<M_1<M_2\), scale any feasible point at \(M_2\) by
\(\sqrt{M_1/M_2}\). Torque becomes \(M_1\), while voltage squared and loss
shrink by \(M_1/M_2\). Hence
\begin{equation}
 P_{\min}(M_1)\leq\frac{M_1}{M_2}P_{\min}(M_2)<P_{\min}(M_2).
 \label{eq:frontier-monotone}
\end{equation}
On each fixed torque-sign branch the attainable frontier is therefore strictly
increasing. Up to the maximum voltage-feasible torque,
\(P_{\min}(M)\) and \(M_{\max}(P)\) are inverse descriptions of the same
efficient boundary. Negative torque follows by selecting directions with
\(T<0\); generator and motor branches need not have identical constraints if
the physical system imposes asymmetric field or inverter limits.

\begin{figure}[tb]
\centering
\begin{tikzpicture}
\begin{axis}[axis main,width=.78\linewidth,height=6cm,
 xlabel={torque $M$},ylabel={minimum loss $P_{\min}$},xmin=0,xmax=5.2,ymin=0,ymax=5.2]
 \addplot[loss curve,domain=0:4.8,samples=100]{.34*x+.12*x^2};
 \addplot[support line] coordinates {(0,3.25) (5.1,3.25)};
 \addplot[support line] coordinates {(3.72,0) (3.72,5)};
 \addplot[only marks,mark=*,mark size=2.2pt,ForestGreen] coordinates {(2.45,1.55)};
 \addplot[only marks,mark=*,mark size=2.2pt,BrickRed] coordinates {(3.72,3.25)};
 \node[annotation] at (axis cs:1.45,1.95){tracking};
 \node[annotation] at (axis cs:4.3,3.75){maximum feasible};
 \node[annotation] at (axis cs:1.2,3.55){$P_{\mathrm{Cu},\max}$};
\end{axis}
\end{tikzpicture}
\caption{Minimum loss and maximum torque are inverse readings of one strictly
increasing efficient frontier on a fixed torque-sign branch.}
\label{fig:frontier}
\end{figure}


This frontier already specifies the controller logic. Compute
\(P_{\mathrm{req}}=P_{\min}(M_{\mathrm{ref}})\). If it is no greater than the
loss ceiling, return the minimizing point. Otherwise return the point on the
same frontier at \(P_{\mathrm{Cu},\max}\), mark the request as limited, and
report its achieved torque.

\paragraph{Synthesis.}
Direction determines efficiency; radius spends resources. Minimum loss fixes
torque and asks for the cheapest direction. Maximum torque fixes resources and
asks which direction turns them into the largest torque. The feasible
operating boundary, not a choice of optimizer, makes them inverse problems.
